
%(BEGIN_QUESTION)
% Copyright 2006, Tony R. Kuphaldt, released under the Creative Commons Attribution License (v 1.0)
% This means you may do almost anything with this work of mine, so long as you give me proper credit

Tyske ubåter i andre verdenskrig brukte svært ren hydrogenperoksid (H$_{2}$O$_{2}$) som monopropellant i torpedofremdriften: når den kommer i kontakt med en katalysator, dekomponerer hydrogenperoksid spontant og danner vann (H$_{2}$O), oksygen (O$_{2}$) og mye varme. Skriv en balansert reaksjonsligning med alle reaktanter og reaksjonsprodukter i riktige proporsjoner.

\vskip 10pt

Bestem også hvor mange {\it mol} vann og hvor mange {\it mol} O$_{2}$ som produseres når nøyaktig tre mol hydrogenperoksid dekomponerer.

\vskip 20pt \vbox{\hrule \hbox{\strut \vrule{} {\bf Forslag til sokratisk diskusjon} \vrule} \hrule}

\begin{itemize}
\item{} Forklar hvordan du kontrollerer at den endelige ligningen er korrekt balansert.
\item{} Bør peroksidets dekomponeringsreaksjon ha en {\it positiv} eller {\it negativ} $\Delta H$-verdi skrevet ved siden av?
\item{} Forklar {\it hvorfor} reaksjonen er eksoterm ut fra reaksjonsligningen: hvor absorberes energi, og hvor frigjøres energi? Hva kan vi si om relative bindingsstyrker for oksygen i H$_{2}$O sammenlignet med oksygen bundet til seg selv?
\item{} Forklar hensikten med å bruke en {\it katalysator} for å starte dekomponeringen av hydrogenperoksid i en slik torpedo.
\item{} Blir katalysatoren forbrukt i dekomponeringsreaksjonen?
\end{itemize}

\underbar{file i00903}
%(END_QUESTION)





%(BEGIN_ANSWER)

\noindent
{\bf Delvis svar:}

$$\hbox{2H}_2\hbox{O}_2 \to \hbox{2H}_2\hbox{O} + \hbox{O}_2$$

%(END_ANSWER)





%(BEGIN_NOTES)

Balansering av reaksjonen med samtidige lineære ligninger:

% No blank lines allowed between lines of an \halign structure!
% I use comments (%) instead, so Tex doesn't choke.

$$\vbox{\offinterlineskip
\halign{\strut
\vrule \quad\hfil # \ \hfil & 
\vrule \quad\hfil # \ \hfil & 
\vrule \quad\hfil # \ \hfil & 
\vrule \quad\hfil # \ \hfil \vrule \cr
\noalign{\hrule}
%
% First row
1 & = & $x$ & $y$ \cr
%
\noalign{\hrule}
%
% Another row
H$_{2}$O$_{2}$ & $\to$ & H$_{2}$O & O$_{2}$ \cr
%
\noalign{\hrule}
} % End of \halign 
}$$ % End of \vbox

% No blank lines allowed between lines of an \halign structure!
% I use comments (%) instead, so Tex doesn't choke.

$$\vbox{\offinterlineskip
\halign{\strut
\vrule \quad\hfil # \ \hfil & 
\vrule \quad\hfil # \ \hfil \vrule \cr
\noalign{\hrule}
%
% First row
{\bf Element} & {\bf Balanse-ligning} \cr
%
\noalign{\hrule}
%
% Another row
Hydrogen & $2 = 2x + 0y$ \cr
%
\noalign{\hrule}
%
% Another row
Oksygen & $2 = 1x + 2y$ \cr
%
\noalign{\hrule}
} % End of \halign 
}$$ % End of \vbox

Fra hydrogenbalansen ($2 = 2x$) ser vi at $x = 1$. Setter vi denne verdien inn i oksygenbalansen ($2 = 1 + 2y$), får vi $y = {1 \over 2}$. For å få heltallskoeffisienter uttrykker vi løsningen som $x = 2$ og $y = 1$, med opprinnelig koeffisient for peroksid økt fra 1 til 2:

$$\hbox{2H}_2\hbox{O}_2 \to \hbox{2H}_2\hbox{O} + \hbox{O}_2$$

\vskip 10pt

Som vi ser av den balanserte ligningen, er molforholdet mellom hydrogenperoksid og vann 1:1, så tre mol peroksid gir tre mol vann. Forholdet mellom peroksid og oksygenmolekyler er 2:1, så tre mol peroksid gir 1,5 mol oksygengass.

\vskip 10pt

Selv om det kreves energi for å frigjøre oksygenatomet fra hydrogenperoksidmolekylet, frigjøres mer energi når frie oksygenatomer binder seg sammen til oksygenmolekyler (O$_{2}$). Derfor er den katalytiske dekomponeringen av hydrogenperoksid {\it eksoterm}.

%INDEX% Chemistry, catalyst
%INDEX% Chemistry, stoichiometry: balancing a chemical equation
%INDEX% Process: torpedo propulsion (hydrogen peroxide)

%(END_NOTES)

